% An invented source paper for the Sextant demo. Nothing in it is real.
\documentclass[10pt]{article}
\usepackage[T1]{fontenc}
\usepackage[margin=2.4cm]{geometry}
\usepackage{lmodern}
\pagestyle{empty}
\begin{document}
\begin{center}
{\small\itshape Journal of Invented Horticulture, vol.~12, pp.~101--109, 2021\par}\vspace{10pt}
{\LARGE\bfseries Moisture Retention in Raised Garden Beds Across Three Growing Seasons\par}\vspace{8pt}
{\large Avery Example and Jordan Sample\par}\vspace{4pt}
{\small Department of Imaginary Studies, Nowhere University\par}
\end{center}
\vspace{6pt}
\noindent\textbf{Abstract.} We followed paired raised and in-ground plots in a shared city garden for three growing seasons and recorded soil moisture weekly at a depth of 15 cm.

\section*{Methods}
Twelve raised beds, each 30 cm high and filled with a loam and compost mix, were paired with twelve in-ground plots of the same size on the same site. Moisture was read with a handheld probe every Monday morning, before any watering, from April to September.

\section*{Results}
Across all three seasons, and in each of the twelve pairs of plots followed from April to September, raised beds retained 23\% more soil moisture than in-ground plots, and the gap was widest in the driest weeks of July and August. In-ground plots also lost more water to run-off after heavy rain, which pooled at the lowest corners of the site.

\section*{Discussion}
The advantage comes mostly from the loose fill: the compacted native soil sheds water that the bed mix absorbs. Beds dried faster at the surface on hot afternoons, so a thin mulch on top is worth adding.
\end{document}
